\documentclass[../../main.tex]{subfiles}
\graphicspath{{../../image/}{../../figure/}}
\begin{document}

\section{解的正则性}\label{section:ch11-解的正则性}

上一节建立的广义特征线性质, 给出了可容许弱解结构与正则性的精确描述。

\begin{theorem}\label{theorem:ch11-11-3-1}
设 $\chi(\cdot)$ 为唯一前向广义特征线, $\xi_-(\cdot),\xi_+(\cdot)$ 为极值后向特征线, 与可容许解 $u$ 关联, 从任意点 $(\bar x,\bar t)\in(-\infty,\infty)\times(0,\infty)$ 出发。则 $(\bar x,\bar t)$ 是函数 $u(x-,t)$ 相对于集合 $\{(x,t):0\leqslant t\leqslant\bar t,\ x\leqslant\xi_-(t)\}$ 或 $\{\bar t<t<\infty,\ x\leqslant\chi(t)\}$ 的连续点, 同时也是函数 $u(x+,t)$ 相对于集合 $\{(x,t):0\leqslant t\leqslant\bar t,\ x\geqslant\xi_+(t)\}$ 或 $\{\bar t<t<\infty,\ x\geqslant\chi(t)\}$ 的连续点。此外, $\chi(\cdot)$ 在 $\bar t$ 处右可微, 且
\begin{align}
\frac{d^+}{dt}\chi(\bar t)=
\begin{cases}
f'(u(\bar x\pm,\bar t)), & \text{若 }u(\bar x+,\bar t)=u(\bar x-,\bar t),\\[2mm]
\dfrac{f(u(\bar x+,\bar t))-f(u(\bar x-,\bar t))}{u(\bar x+,\bar t)-u(\bar x-,\bar t)}, & \text{若 }u(\bar x+,\bar t)<u(\bar x-,\bar t).
\end{cases}\label{align:ch11-11-3-1}
\end{align}
\end{theorem}

\begin{proof}
取集合 $\{(x,t):0\leqslant t\leqslant\bar t,\ x\leqslant\xi_-(t)\}$ 或 $\{\bar t<t<\infty,\ x\leqslant\chi(t)\}$ 中任意序列 $\{(x_n,t_n)\}$, 当 $n\to\infty$ 时收敛于 $(\bar x,\bar t)$。设 $\xi_n(\cdot)$ 表示从 $(x_n,t_n)$ 出发的极小后向特征线。显然, 对 $t\leqslant\bar t$, $\xi_n(t)\leqslant\xi_-(t)$。于是当 $n\to\infty$ 时, $\{\xi_n(\cdot)\}$ 从下方收敛到 $\xi_-(\cdot)$。因此, $\{u(x_n-,t_n)\}$ 收敛到 $u(\bar x-,\bar t)$。

类似地, 对集合 $\{(x,t):0\leqslant t<\bar t,\ x\geqslant\xi_+(t)\}$ 或 $\{\bar t<t<\infty,\ x\geqslant\chi(t)\}$ 中任意收敛于 $(\bar x,\bar t)$ 的序列 $\{(x_n,t_n)\}$, 序列 $\{u(x_n+,t_n)\}$ 收敛到 $u(\bar x+,\bar t)$。

对 $\varepsilon>0$,
\begin{align}
\frac{1}{\varepsilon}[\chi(\bar t+\varepsilon)-\chi(\bar t)]=\frac{1}{\varepsilon}\int_{\bar t}^{\bar t+\varepsilon}\dot{\chi}(t)\,\mathrm{d}t,\label{align:ch11-11-3-2}
\end{align}
其中 $\dot{\chi}(t)$ 通过\eqref{align:ch11-11-1-3}确定, 取 $\xi\equiv\chi$。如上所示, $\dot{\chi}(t)$ 在 $\bar t$ 处右连续, 故在\eqref{align:ch11-11-3-2}中令 $\varepsilon\downarrow0$ 即得\eqref{align:ch11-11-3-1}。证毕。
\end{proof}

上述定理有如下推论。

\begin{theorem}\label{theorem:ch11-11-3-2}
设 $u$ 为可容许解, 并假设对某个 $(\bar x,\bar t)\in(-\infty,\infty)\times(0,\infty)$, $u(\bar x+,\bar t)=u(\bar x-,\bar t)$。则 $(\bar x,\bar t)$ 是 $u$ 的连续点。存在唯一广义特征线 $\chi(\cdot)$, 与 $u$ 关联, 定义在 $[0,\infty)$ 上, 经过 $(\bar x,\bar t)$。此外, $\chi(\cdot)$ 在 $\bar t$ 处可微, 且 $\dot{\chi}(\bar t)=f'(u(\bar x\pm,\bar t))$。
\end{theorem}

接下来关注间断点。

\begin{theorem}\label{theorem:ch11-11-3-3}
设 $u$ 为可容许解, 并假设对某个 $(\bar x,\bar t)\in(-\infty,\infty)\times(0,\infty)$, $u(\bar x+,\bar t)<u(\bar x-,\bar t)$。当极值后向特征线 $\xi_-(\cdot),\xi_+(\cdot)$ 是从 $(\bar x,\bar t)$ 出发且在 $(0,\bar t)$ 上无激波的仅有的后向广义特征线时, $(\bar x,\bar t)$ 是 $u$ 的跳跃间断点, 意义如下: 存在一条广义特征线 $\chi(\cdot)$, 与 $u$ 关联, 定义在 $[0,\infty)$ 上并经过 $(\bar x,\bar t)$, 使得 $(\bar x,\bar t)$ 是函数 $u(x-,t)$ 相对于 $\{(x,t):0<t<\infty,\ x\leqslant\chi(t)\}$ 的连续点, 同时也是函数 $u(x+,t)$ 相对于 $\{(x,t):0<t<\infty,\ x\geqslant\chi(t)\}$ 的连续点。此外, $\chi(\cdot)$ 在 $\bar t$ 处可微, 且
\begin{align}
\dot{\chi}(\bar t)=\frac{f(u(\bar x+,\bar t))-f(u(\bar x-,\bar t))}{u(\bar x+,\bar t)-u(\bar x-,\bar t)}.\label{align:ch11-11-3-3}
\end{align}
\end{theorem}

\begin{proof}
在 $x$ 轴上区间 $(\xi_-(0),\xi_+(0))$ 内任取一点, 用特征线 $\chi(\cdot)$ 将其与 $(\bar x,\bar t)$ 连接。将 $\chi(\cdot)$ 从 $[\bar t,\infty)$ 延拓为从 $(\bar x,\bar t)$ 出发的唯一前向特征线。

取集合 $\{(x,t):0<t<\infty,\ x\leqslant\chi(t)\}$ 中任意序列 $\{(x_n,t_n)\}$ 收敛于 $(\bar x,\bar t)$。设 $\xi_n(\cdot)$ 表示从 $(x_n,t_n)$ 出发的极小后向特征线。当 $n\to\infty$ 时, $\{\xi_n(\cdot)\}$ 或其子列将收敛到某条从 $(\bar x,\bar t)$ 出发的后向特征线, 它是直线且无激波。由于 $\xi_n(t)\leqslant\chi(t)$, 这蕴含 $\{\xi_n(\cdot)\}$ 必收敛到 $\xi_-(\cdot)$。因此, 当 $n\to\infty$ 时, $\{u(x_n-,t_n)\}$ 收敛到 $u(\bar x-,\bar t)$。

类似地, 对 $\{(x,t):0<t<\infty,\ x\geqslant\chi(t)\}$ 中任意收敛于 $(\bar x,\bar t)$ 的序列 $\{(x_n,t_n)\}$, 序列 $\{u(x_n+,t_n)\}$ 收敛到 $u(\bar x+,\bar t)$。

为验证\eqref{align:ch11-11-3-3}, 我们从\eqref{align:ch11-11-3-2}重新出发, 此时 $\varepsilon$ 可正可负。如上所示, $\bar t$ 是 $\dot{\chi}(t)$ 的连续点, 故令 $\varepsilon\to0$ 即得\eqref{align:ch11-11-3-3}。证毕。
\end{proof}

\begin{theorem}\label{theorem:ch11-11-3-4}
任意可容许解 $u$ 的不规则点集是可数的。点 $(\bar x,\bar t)\in(-\infty,\infty)\times(0,\infty)$ 是不规则点当且仅当 $u(\bar x+,\bar t)<u(\bar x-,\bar t)$, 并且除极值后向特征线 $\xi_-(\cdot),\xi_+(\cdot)$ 外, 至少还有另一条不同的、与 $u$ 关联的后向特征线 $\xi(\cdot)$ 从 $(\bar x,\bar t)$ 出发, 且在 $(0,\bar t)$ 上无激波。不规则点由激波碰撞和/或中心压缩波聚焦产生。
\end{theorem}

\begin{proof}
必要性由\cref{theorem:ch11-11-3-2}与\cref{theorem:ch11-11-3-3}给出。为证充分性, 考虑区间 $[\xi_-(0),\xi_+(0)]$ 的子集 $\mathscr{X}$: 对 $x\in\mathscr{X}$, 连接 $(x,0)$ 与 $(\bar x,\bar t)$ 的直线段是与 $u$ 关联的特征线, 且在 $(0,\bar t)$ 上无激波。

当 $\mathscr{X}\equiv[\xi_-(0),\xi_+(0)]$ 时, $(\bar x,\bar t)$ 是中心压缩波的中心, 定理断言显然成立。一般地, $\mathscr{X}$ 是 $[\xi_-(0),\xi_+(0)]$ 的闭真子集, 至少包含 $\xi_-(0),\xi(0)$ 与 $\xi_+(0)$ 三个点。$\mathscr{X}$ 相对于 $[\xi_-(0),\xi_+(0)]$ 的余集将是(至多)可数个不相交开区间的并。设 $(\alpha_-,\alpha_+)$ 为其中一个区间, 不妨它含于 $(\xi_-(0),\xi(0))$。连接 $(\alpha_-,0)$ 与 $(\alpha_+,0)$ 和 $(\bar x,\bar t)$ 的直线段是无激波特征线 $\zeta_-(\cdot)$ 与 $\zeta_+(\cdot)$, 沿其上 $u$ 为常数, 记为 $u_-$ 与 $u_+$。必有 $u(\bar x-,\bar t)\geqslant u_->u_+>u(\bar x+,\bar t)$。考虑连接 $(\alpha_-,\alpha_+)$ 中某点与 $(\bar x,\bar t)$ 的特征线 $\chi(\cdot)$。则 $\zeta_-(t)<\chi(t)<\zeta_+(t)$, $0\leqslant t<\bar t$。取集合 $\{(x,t):0\leqslant t<\bar t,\ \zeta_-(t)\leqslant x\leqslant\chi(t)\}$ 中任意收敛于 $(\bar x,\bar t)$ 的序列 $\{(x_n,t_n)\}$。若 $\xi_n(\cdot)$ 表示从 $(x_n,t_n)$ 出发的极小后向特征线, 序列 $\{\xi_n(\cdot)\}$ 必收敛到 $\zeta_-(\cdot)$。特别地, $u(x_n-,t_n)\to u_-$, $n\to\infty$。类似地, 若 $\{(x_n,t_n)\}$ 是集合 $\{(x,t):0\leqslant t<\bar t,\ \chi(t)\leqslant x\leqslant\zeta_+(t)\}$ 中任意收敛于 $(\bar x,\bar t)$ 的序列, 则 $u(x_n+,t_n)\to u_+$, $n\to\infty$。于是在 $\bar t$ 附近, $\chi(\cdot)$ 是激波, 在 $\bar t$ 处左可微, 且
\begin{align}
\frac{d^-}{dt}\chi(\bar t)=\frac{f(u_+)-f(u_-)}{u_+-u_-}.\label{align:ch11-11-3-4}
\end{align}
由于 $f'(u_-)>\dfrac{d^-}{dt}\chi(\bar t)>f'(u_+)$, 我们断定 $(\bar x,\bar t)$ 是 $u$ 的不规则点。

于是我们证明了 $(\bar x,\bar t)$ 是激波碰撞点, 每个开区间对应一次碰撞; 当 $\mathscr{X}$ 的测度为正时, 还有中心压缩波。

对固定正数 $\varepsilon$, 考虑上述不规则点 $(\bar x,\bar t)$, 附加性质 $\xi_+(0)-\xi(0)>\varepsilon$, $\xi(0)-\xi_-(0)>\varepsilon$。易见在上半平面的任意有界子集内, 这样的点至多可容纳有限个。这进而推出任意可容许解的不规则点集是(至多)可数的。证毕。
\end{proof}

真正非线性的影响, 反映在特征线性质上, 要么通过稀疏波光滑化解, 要么通过压缩形成跳跃间断。这种极化影响阻碍了具有"中间"正则性的解的存在, 其体现见下文\cref{theorem:ch11-11-3-5}、\cref{theorem:ch11-11-3-6}与\cref{theorem:ch11-11-3-10}。

首先, 每个可容许 $BV$ 解必定是\cref{chapter:ch01-平衡律}意义下的特殊有界变差函数。


\begin{theorem}\label{theorem:ch11-11-3-5}
存在一个(至多)可数集 $\mathscr{T}\subset[0,\infty)$, 使得对任意 $t\in[0,\infty)\setminus\mathscr{T}$, $u(\cdot,t)\in SBV_{\mathrm{loc}}(-\infty,\infty)$。此外, $u\in SBV_{\mathrm{loc}}((-\infty,\infty)\times[0,\infty))$。
\end{theorem}

上述命题的证明见\cref{chapter:ch11-一维标量守恒律}所引文献。这一现象的粗略解释如下。由\cref{theorem:ch11-11-2-2}, 对固定 $t>0$, 解 $u$ 关于 $x$ 可下降, 但不能快速上升。此外, 由特征线性质, 连续而不突变的快速下降只能在压缩波破裂前不久(特别是在中心压缩波聚焦的临界点附近)出现。$u(\cdot,t)$ 的 Cantor 部分的发展, 可能由无穷多个压缩波几乎同时破裂所诱发, 这是罕见的、非一般的事件。


接下来我们将看到, 连续性自动升级为 Lipschitz 连续。

\begin{theorem}\label{theorem:ch11-11-3-6}
假设可容许解 $u$ 的连续点集 $\mathscr{C}$ 具有非空内部 $\mathscr{C}^0$。则 $u$ 在 $\mathscr{C}^0$ 上局部 Lipschitz 连续。
\end{theorem}

\begin{proof}
固定任意点 $(\bar x,\bar t)\in\mathscr{C}^0$, 并假设以 $(\bar x,\bar t)$ 为心、半径为 $r$ 的圆盘 $\mathscr{B}_r$ 含于 $\mathscr{C}^0$。考虑距 $(\bar x,\bar t)$ 距离为 $\rho<r$ 的任意点 $(x,t)$。与 $u$ 关联、经过 $(\bar x,\bar t)$ 与 $(x,t)$ 的(唯一)特征线分别是斜率为 $f'(u(\bar x,\bar t))$ 与 $f'(u(x,t))$ 的直线, 它们不能在圆盘 $\mathscr{B}_r$ 内相交。初等三角估计进而推出 $|f'(u(x,t))-f'(u(\bar x,\bar t))|$ 不超过 $c\rho/r$, 其中 $c$ 是 $1+f'(u)^2$ 在 $\mathscr{B}_r$ 上的上界。因此, 若 $a>0$ 是 $f''(u)$ 在 $\mathscr{B}_r$ 上的下界, 则 $|u(x,t)-u(\bar x,\bar t)|\leqslant\dfrac{c}{ar}\rho$。证毕。
\end{proof}

读者应注意, 已构造出连续点集内部为空的可容许解。

现在研究初值光滑的可容许解的正则性。下文假设 $f$ 是 $C^{k+1}$ 的, $u$ 是初值 $u_0\in C^k$ 的可容许解, 其中 $k\in\{1,2,\cdots,\infty\}$。

对 $(x,t)\in(-\infty,\infty)\times(0,\infty)$, 设 $y_-(x,t)$ 与 $y_+(x,t)$ 表示从点 $(x,t)$ 出发的、与 $u$ 关联的极小与极大后向特征线在 $x$ 轴上的截距。特别地,
\begin{align}
x=y_-(x,t)+tf'(u_0(y_-(x,t)))=y_+(x,t)+tf'(u_0(y_+(x,t))),\label{align:ch11-11-3-5}\\
u(x-,t)=u_0(y_-(x,t)),\qquad u(x+,t)=u_0(y_+(x,t)).\label{align:ch11-11-3-6}
\end{align}

对固定 $t>0$, $y_-(\cdot,t)$ 与 $y_+(\cdot,t)$ 均单调不减, 前者左连续而后者右连续。因此, 对所有 $(x,t)\in(-\infty,\infty)\times(0,\infty)$,
\begin{align}
1+t\frac{d}{dy}f'(u_0(y))\geqslant0,\quad y=y_\pm(x,t)\label{align:ch11-11-3-7}
\end{align}
成立。

$u$ 的任意连续点 $(\bar x,\bar t)\in(-\infty,\infty)\times(0,\infty)$ 也必是 $y_\pm(x,t)$ 的连续点, 且 $y_-(\bar x,\bar t)=y_+(\bar x,\bar t)$。因此, 由\eqref{align:ch11-11-3-5}、\eqref{align:ch11-11-3-6}与隐函数定理推出

\begin{theorem}\label{theorem:ch11-11-3-7}
若 $(\bar x,\bar t)\in(-\infty,\infty)\times(0,\infty)$ 是 $u$ 的连续点, 且
\begin{align}
1+\bar t\frac{d}{dy}f'(u_0(y))>0,\quad y=y_\pm(\bar x,\bar t),\label{align:ch11-11-3-8}
\end{align}
则 $u$ 在 $(\bar x,\bar t)$ 的邻域内是 $C^k$ 的。
\end{theorem}

参照\cref{theorem:ch11-11-3-3}, 若 $(\bar x,\bar t)$ 是 $u$ 的跳跃间断点, 则 $(\bar x,\bar t)$ 分别是 $y_-(x,t)$ 相对于点集 $\{(x,t):0<t<\infty,\ x\leqslant\chi(t)\}$ 和 $y_+(x,t)$ 相对于 $\{(x,t):0<t<\infty,\ x\geqslant\chi(t)\}$ 的连续点。于是隐函数定理连同\eqref{align:ch11-11-3-5}与\eqref{align:ch11-11-3-6}给出

\begin{theorem}\label{theorem:ch11-11-3-8}
若在 $u$ 的跳跃间断点 $(\bar x,\bar t)\in(-\infty,\infty)\times(0,\infty)$ 处\eqref{align:ch11-11-3-8}成立, 则在 $(\bar x,\bar t)$ 的邻域内, 经过 $(\bar x,\bar t)$ 的激波 $\chi(\cdot)$ 是 $C^{k+1}$ 的, 而 $u$ 在 $\chi(\cdot)$ 图像两侧分别是 $C^k$ 的。
\end{theorem}

接下来考虑\cref{definition:ch11-11-6}引入的激波产生点。

\begin{theorem}\label{theorem:ch11-11-3-9}
若 $(\bar x,\bar t)\in(-\infty,\infty)\times(0,\infty)$ 是激波产生点, 则
\begin{align}
1+\bar t\frac{d}{dy}f'(u_0(y))=0,\quad y_-(\bar x,\bar t)\leqslant y\leqslant y_+(\bar x,\bar t).\label{align:ch11-11-3-9}
\end{align}
此外, 当 $k\geqslant2$ 时,
\begin{align}
\frac{d^2}{dy^2}f'(u_0(y))=0,\quad y_-(\bar x,\bar t)\leqslant y\leqslant y_+(\bar x,\bar t).\label{align:ch11-11-3-10}
\end{align}
\end{theorem}

\begin{proof}
回忆激波产生点有两类: 连续点情形, 此时 $y_-(\bar x,\bar t)=y_+(\bar x,\bar t)$; 中心压缩波中心情形, 此时 $y_-(\bar x,\bar t)<y_+(\bar x,\bar t)$。当 $(\bar x,\bar t)$ 是连续点时, \eqref{align:ch11-11-3-9} 是\eqref{align:ch11-11-3-7}与\cref{theorem:ch11-11-3-7}的推论。当 $(\bar x,\bar t)$ 是压缩波中心时, 对任意 $y\in[y_-(\bar x,\bar t),y_+(\bar x,\bar t)]$, $\bar x=y+\bar t f'(u_0(y))$, 这蕴含\eqref{align:ch11-11-3-9}。

当 $y_-(\bar x,\bar t)<y_+(\bar x,\bar t)$ 时, 对 $y$ 微分\eqref{align:ch11-11-3-9}得\eqref{align:ch11-11-3-10}。为在 $(\bar x,\bar t)$ 为连续点的情形建立\eqref{align:ch11-11-3-10}, 取任意从下方(或上方)收敛到 $\bar x$ 的序列 $\{x_n\}$。则 $\{y_-(x_n,\bar t)\}$ 将从下方(或上方)趋近 $y_\pm(\bar x,\bar t)$。由于对 $y=y_-(x_n,\bar t)$, $1+\bar t\dfrac{d}{dy}f'(u_0(y))\geqslant0$; 这连同\eqref{align:ch11-11-3-9}推出 $y_\pm(\bar x,\bar t)$ 是 $\dfrac{d}{dy}f'(u_0(y))$ 的临界点。证毕。
\end{proof}

对 $k\geqslant3$, $C^k$ 函数 $u_0$ 中满足 $\dfrac{d}{dy}f'(u_0(y))$ 在有界区间内有无穷多临界点的函数集是第一范畴的。因此, 一般地, 初值 $u_0\in C^k$, $k\geqslant3$, 诱导出的解具有局部有限的激波产生点集, 从而具有局部有限的激波集。换言之,

\begin{theorem}\label{theorem:ch11-11-3-10}
一般地, 初值属于 $C^k$, $k\geqslant3$ 的\eqref{align:ch11-11-1-1}的可容许解是分段 $C^k$ 光滑函数, 且不含中心压缩波。特别地, 初值解析的解总是分段解析的。
\end{theorem}

\end{document}
